\chapter{Геометрический родитель взаимно обратного спектрального хода}
\label{ch:v10-k43-reciprocal-spectral-operator-growth-parent-origin}

\section{Проверяемая гипотеза}

Предыдущий гейт нашёл внутри $\mathcal K_{43}$ каноническую пару
$\mathcal W_Y=\operatorname{span}\{|Y\rangle,|0\rangle\}$, но оператор с
собственными значениями $e^{-\zeta}$ и $e^{\zeta}$ был задан вручную.
Настоящий гейт проверяет, можно ли вывести эту зависимость как единственную
нулевую траекторию положительного геометрического родителя, используя только
два унаследованных проектора и стрелку притока.

\section{Канонический генератор роста}

Обозначим проекторы на гиперзарядовую и вакуумную линии через $P_Y$ и $P_0$.
Разность
\begin{equation}
 Q_X:=P_0-P_Y
 \label{eq:v10-growth-grading}
\end{equation}
обладает точными свойствами
\begin{equation}
 \Tr Q_X=0,
 \qquad Q_X^2=P_0+P_Y,
 \qquad \rank Q_X=2,
 \label{eq:v10-growth-grading-identities}
\end{equation}
а её спектр на $\mathcal K_{43}$ равен
\begin{equation}
 \Spec Q_X=\{-1^{(1)},0^{(41)},1^{(1)}\}.
 \label{eq:v10-growth-grading-spectrum}
\end{equation}

Безследовость является условием сохранения определителя. На диагональном
двумерном носителе она оставляет одномерное пространство генераторов.
Нормировка $\Tr Q_X^2=2$ оставляет только два знака $\pm Q_X$. Стрелка
притока выбирает $Q_X$, поскольку её согласованности равны
\begin{equation}
 \Tr(Q_XQ_X)=2,
 \qquad
 \Tr(Q_X(-Q_X))=-2.
 \label{eq:v10-growth-arrow-sign-selection}
\end{equation}
Таким образом, веса $(-1,+1)$ не вводятся как два независимых коэффициента.
Они следуют из сохранения определителя, единичной нормировки координаты
$\zeta$ и ориентации роста.

\section{Положительный родитель траектории}

Для положительного операторного пути $K(\zeta)$ определим симметричную
ковариантную производную
\begin{equation}
 \nabla_XK
 :=\frac{dK}{d\zeta}-\frac12\{Q_X,K\}.
 \label{eq:v10-growth-covariant-derivative}
\end{equation}
Минимальный родитель пути имеет вид
\begin{equation}
 \mathcal S_X[K]
 =\frac12\|K(0)-I_{43}\|_{\rm HS}^2
 +\frac12\int_0^Z\|\nabla_XK(\zeta)\|_{\rm HS}^2\,d\zeta.
 \label{eq:v10-growth-path-parent}
\end{equation}
Он неотрицателен. Нулевое значение требует одновременно
\begin{equation}
 K(0)=I_{43},
 \qquad
 \nabla_XK=0.
 \label{eq:v10-growth-parent-zero-equations}
\end{equation}
Единственное решение этой линейной начальной задачи равно
\begin{equation}
 K_X(\zeta)=e^{\zeta Q_X}
 =I_{43}+(e^{-\zeta}-1)P_Y+(e^{\zeta}-1)P_0.
 \label{eq:v10-growth-parent-solution}
\end{equation}
Поскольку $\Tr Q_X=0$,
\begin{equation}
 \det K_X(\zeta)=1.
 \label{eq:v10-growth-parent-determinant}
\end{equation}

Здесь экспоненциальная кривая является следствием дифференциального
уравнения и начального условия, а не целью, заранее записанной внутри
потенциала.

\section{Локальная строгая проверка родителя}

На диагональном начальном струйном наборе
$(r_Y,r_0,v_Y,v_0)$ конечномерный представитель равен
\begin{equation}
 \mathcal L_X
 =\frac12(v_Y+r_Y)^2+\frac12(v_0-r_0)^2
 +\frac12(r_Y-1)^2+\frac12(r_0-1)^2.
 \label{eq:v10-growth-parent-jet}
\end{equation}
Точка
\begin{equation}
 (r_Y,r_0,v_Y,v_0)=(1,1,-1,1)
 \label{eq:v10-growth-parent-jet-point}
\end{equation}
имеет нулевой градиент и нулевое значение. Её гессиан имеет
\begin{equation}
 \rank\operatorname{Hess}\mathcal L_X=4,
 \qquad
 \det\operatorname{Hess}\mathcal L_X=1.
 \label{eq:v10-growth-parent-jet-hessian}
\end{equation}
Следовательно, начальное условие и обе ориентированные скорости выбраны
строго, без плоской струйной моды.

\section{Возвращение самоэнергии}

Сжатие решения на $\mathcal W_Y$ даёт
\begin{equation}
 E_Y^*K_X(\zeta)E_Y
 =\operatorname{diag}(e^{-\zeta},e^{\zeta}).
 \label{eq:v10-growth-parent-compression}
\end{equation}
Поэтому входящая самоэнергия теперь следует из родительской траектории:
\begin{equation}
 \Sigma_Y(\zeta)
 =\langle Y|K_X^{-1}(\zeta)|Y\rangle=e^{\zeta},
 \qquad
 \frac{d\Sigma_Y}{d\zeta}=\Sigma_Y.
 \label{eq:v10-growth-parent-self-energy}
\end{equation}

\begin{theorem}[Структурное происхождение взаимно обратного хода]
Безследовый нормированный генератор $Q_X=P_0-P_Y$, выбранный стрелкой
притока, и неотрицательный родитель \eqref{eq:v10-growth-path-parent}
однозначно выбирают $K_X(\zeta)=e^{\zeta Q_X}$. Тем самым взаимно обратный
спектр и ненулевая интенсивная самоэнергия получены без подстановки готовой
экспоненциальной кривой в потенциал.
\end{theorem}

\begin{nogocriterion}[Геометрическая координата ещё не является секундой]
Родитель определяет зависимость от безразмерной координаты
$\zeta=\frac13\log(N/N_0)$, но не задаёт вероятность перехода $N\to N+1$
и не переводит один шаг истории в физическое время или энергию. Поэтому
структурное происхождение спектрального закона закрыто, а абсолютный
размерный масштаб и физическая космологическая постоянная ещё не выведены.
\end{nogocriterion}

\begin{protocol}[Статус геометрического родителя]\begin{enumerate}
\item безследовый генератор роста: \textbf{выведен};
\item выбор знака стрелкой притока: \textbf{$2>-2$};
\item положительный родитель пути: \textbf{построен};
\item единственная нулевая траектория: \textbf{$e^{\zeta Q_X}$};
\item архитектура геометрического родителя: \textbf{$9/9$};
\item структурное происхождение спектрального закона: \textbf{$4/4$};
\item нормированная мера рождения ячеек: \textbf{$0/1$};
\item физическая временно-энергетическая калибровка: \textbf{$0/1$};
\item абсолютный масштаб: \textbf{не выведен};
\item следующий гейт: \textbf{нормированная мера переходов рождения ячеек
и происхождение темпа роста}.
\end{enumerate}\end{protocol}

Машинный сертификат сохранён в
\path{s2t_v10_k43_reciprocal_spectral_operator_growth_parent_origin_gate_results.json}.